\documentclass[a4paper,12pt]{article}
\usepackage[utf8]{inputenc}
\usepackage[ngerman]{babel}
\usepackage{amsmath}
\usepackage{geometry}
\geometry{margin=2.5cm}
\usepackage{booktabs}
\usepackage{parskip}
\usepackage{tikz}
\usepackage{href-ul}
\hypersetup{
	colorlinks=true,
	linkcolor=blue,
	urlcolor=blue}
\date{}

\begin{document}
	
	\begin{tikzpicture}(10,3)
		\draw[ultra thick](2,0) --(10,0) -- (10,1.3) --(2,1.3) -- (2,0);
		\draw[fill=black](2,0)-- (0,.6) -- (2,.6) -- (2,0);
		\node at (6,.6) {\large solution sheet from \href{https://www.okuyakl.de}{www.okuyakl.de}};
	\end{tikzpicture}
	
	\vspace{0.5 cm}
	
	\section*{1. The Hungry Hamster on the Moon}
	
	\begin{itemize}
		\item a) \( m = 0.150\,\text{kg} \), \( g_\text{Earth} = 9.81\,\frac{\text{m}}{\text{s}^2 }\)
			
			\[
			F = 0.150 \cdot 9.81 = \boxed{1.47\,\text{N}}
			\]
			
			\item b) \( g_\text{Moon} = 1.6 \)
			
			\[
			F = 0.150 \cdot 1.6 = \boxed{0.24\,\text{N}}
			\]
			
			\item c) Difference in percent:
			
			\[ 
			\frac{1{,}47 - 0{,}24}{1{,}47} \cdot 100 \approx \boxed{83{,}7\,\%}
			\]
		\end{itemize}
		
		\hrulefill
		
		\section*{2. The elephant's yoga pose}
		
		\begin{itemize}
			\item a) \( F = 3000 \cdot 9{,}81 = \boxed{29430\,\text{N}} \)
			
			\item b) Rearrange: \( m = \frac{F}{g} = \frac{33\,000}{10} = \boxed{3300\,\text{kg}} \)
		\end{itemize}
		
		\hrulefill
		
		\section*{3. The helium balloon in the sauna}
		
		\begin{itemize}
			\item \( m = 0{.}050\,\text{kg} \), \( g = 9{.}81 \)
			
			\[
			F = 0.050 \cdot 9.81 = \boxed{0.49\,\text{N}}
			\]
		\end{itemize}
		
		\hrulefill
		
		\section*{4. The intergalactic luggage scale}
		
		\begin{itemize}
			\item \( m = 25\,\text{kg} \)
			
			\begin{itemize}
				\item Earth: \( 25 \cdot 9{.}81 = \boxed{245{.}25\,\text{N}} \)
				\item Mars: \( 25 \cdot 3{,}71 = \boxed{92{,}75\,\text{N}} \)
				\item Jupiter: \( 25 \cdot 24{.}79 = \boxed{619{.}75\,\text{N}} \)
				\item Uranus: \( 25 \cdot 8{,}69 = \boxed{217{}25\,\text{N}} \)
			\end{itemize}
			
			\item b) Heaviest: \boxed{Jupiter}
			
		\end{itemize}
		
		\hrulefill
		
		\section*{5. The Horseshoe in Love}
		
		\begin{itemize}
			\item a) \( m = 2\,\text{kg} \), \( g = 9{}81 \), so:
			
			\[
			F = 2 \cdot 9{}81 = \boxed{19{}62\,\text{N}}
			\]
			
			\item b) The magnetic force must be at least \boxed{19{}62\,\text{N}}, otherwise the iron will fall.
			
			\item c) Love is measurable when it has an attractive force. In this case: \textbf{19.62 N romantic tension}.
		\end{itemize}
		
		\begin{center}
			\includegraphics[width=7 cm]{../../viecher/eendcomic.pdf}
			
			Click here to return to the \href{https://www.okuyakl.de/phys/p7fmgL058/ae058.pdf}{worksheet}
		\end{center}
	\end{document}